% ============================================================
% NOTE (draft status, 2026-07): レポート調ドラフト.
% ============================================================

\section{Conclusions}
\label{sec:conclusions}

揚水発電上池を想定した浅い矩形 basin(内寸 \SI{2.952}{m} $\times$ \SI{1.962}{m})において, 単一の可逆開口(幅 \SI{0.11}{m}, 高さ $a=\SI{0.055}{m}$)を介した inflow / outflow の表面流を time-resolved 表面 PIV で計測し, 相対水深 $h/a = 0.55$--$1.82$ にわたる表面流レジームを 24 run(mode 2 $\times$ 流量 3 $\times$ 反復 4, 全て静水開始)で整理した. inflow の 12 run は 4 日 $\times$ 3 流量水準の block 計画とし, 流量の効果を日ごとの条件から分離できるようにした. 主要な結論は以下のとおりである.

\begin{enumerate}
\item 同一開口を共有していても, inflow と outflow の表面流レジームは決定的に非対称である. inflow は全水深帯で「偏向ジェット + 非対称循環」($I_{\mathrm{asym}} \approx 0.26$, $I_{\mathrm{rot}} = 0.06$--$0.12$)を示すのに対し, outflow は全水深帯で回転をもたない「対称収束シンク」($I_{\mathrm{asym}} \approx 0.08$, $I_{\mathrm{rot}} \le 0.007$)にとどまり, 表面運動強度 $E$ は約 3 桁小さい. この差は流入(運動量をもつジェット)と流出(分布した吸込み)の力学的非対称に帰着し, 浅い幾何($A_{\mathrm{plan}}/(ab) \sim 10^3$)がこれを増幅する.

\item inflow の循環の\textbf{向き}は流量によって選択される. 深水帯($h/a>1.09$)の符号付き循環指標でみると, low は 4/4 が反時計回り, high は 4/4 が時計回りで例外がなく, medium のみが 3 対 1 に割れた. inflow は 4 日 $\times$ 3 水準の block 計画で実施しており, 同一日の中で low と high の符号が必ず逆になることから, この選択は日ごとの設営条件ではなく流量に帰属する. medium で逆符号を示した run の実測流量は同水準の他 run と \SI{1.4}{\percent} しか違わず, 同一条件下で二つの分枝がともに実現しうる(双安定). さらに分枝の固定の強さ $R_{\mathrm{lock}}$ は流量とともに単調に増大し(水準中央値 3.6, 8.7, 11.8), 高流量では 12 run 中 10 run で符号保持率が 1.000 に達した. 流量は, 分枝の選択と固定の強さの双方を支配する制御パラメータとして働く.

\item 相対水深 $h/a \approx 1$ は, 流況トポロジーの切り替えではなく, 非定常性と表面影響度の転換帯として機能する. inflow の帯別非定常指標 $I_{\mathrm{unst}}$ は $h < a$ で 0.35--0.40, $h > a$ で 0.03--0.05 と約 1 桁の急減を示し, この遷移は全流量水準で共通であった. 低流速域率 $\phi_{\mathrm{lv}}$($u_{\mathrm{th}} = 0.2\,U_p$)も $h/a$ とともに単調増加した.

\item outflow の表面流は, 時間平均が微弱である一方, 相対変動は大きい. 解像された条件では $I_{\mathrm{unst}} = 0.76$ と inflow 最大値の約 2 倍に達し, 「平均場のみでは outflow を特徴付けられない」ことを basin スケールで定量化した. また, 全 outflow 条件で表面全域が $0.2\,U_p$ 未満($\phi_{\mathrm{lv}} = 1$)であり, port の駆動が自由表面へほとんど到達しないこと自体が浅い可逆 basin の outflow の本質的特徴である.

\item 符号付き循環指標 $I_{\mathrm{circ}}$ は両 mode で 0 近傍となるが, 由来は正反対(inflow: 逆符号循環の相殺, outflow: 回転の不在)であり, 回転活動度 $I_{\mathrm{rot}}$ との組で初めて区別できる. 指標設計上の教訓として, 閉じた basin の循環評価には符号付き・符号なしの併用が必要である.
\end{enumerate}

今後の課題として, (i) 履歴効果(inflow 直後の outflow)の系統的検証, (ii) outflow 非定常性の分解能を確保した計測設計, (iii) 分枝の向きを決めている非対称の同定(給水系の非対称性を意図的に振る実験), (iv) 底面せん断・堆砂分布との対応付け, が挙げられる.

\appendix
\section{Appendix}
\subsection{Additional QC plots and sensitivity checks}
